% !TEX program = xelatex
% 通用中文学习笔记模板（study-note-style-v1）
% 编译：先运行 scripts/compile_study_note.py --tex 本文件；仅在目录、书签或引用未稳定时重编译。
% 本文件与 study-note-style.sty 必须位于同一目录。
\documentclass[12pt,a4paper,fontset=none]{ctexart}
\usepackage{study-note-style}

\begin{document}

% 紧凑标题区
\begin{center}
  {\Huge\bfseries 课程名称 \\[0.5ex]}
  {\LARGE 副标题或资料范围 \\[1ex]}
  {\small 教师姓名 · 院系 · 学期}

  \vspace{0.4em}
  {\small 标注说明：\PriorityMust\ 必背\;|\;\PriorityImportant\ 重点\;|\;\PriorityKnow\ 了解}

  \vspace{0.4em}
  {\small\textbf{阅读顺序：}核心概念 $\rightarrow$ 推导与例题 $\rightarrow$ 公式速查 $\rightarrow$ 易错复盘}
\end{center}
\StudyDivider

% 目录保留点击能力，但 study-note-style 使用 hidelinks，不显示红框或彩色链接。
\tableofcontents
\clearpage

% 此单章节模板展示“核心知识关系导图”。1--3 个正文主章节用核心概念或机制做节点；
% 4 章及以上改为“各章节关系导图”，节点对应实际章节。节点和连线须按资料替换，
% 且正文须解释每个节点与关系；若无可靠关系可画，删去导图并写明 % study-map-omitted: 原因。
\section*{核心知识关系导图}
\addcontentsline{toc}{section}{核心知识关系导图}
\begin{center}
  {\small\textbf{说明：}本示意图以正文中的概念、条件和推导为节点；生成实际笔记时须依据资料重写每条关系。}
\end{center}
\vspace{0.3em}

\begin{center}
\resizebox{0.98\textwidth}{!}{%
\begin{tikzpicture}[study map base]
  \node[study map foundation] (concept) {核心概念\\术语名称};
  \node[study map core,below left=1.5cm and 2.5cm of concept] (condition) {成立条件\\适用边界};
  \node[study map core,below right=1.5cm and 2.5cm of concept] (derivation) {推导与应用\\公式速记};
  \node[study map core,below=2cm of concept] (result) {关键结论\\公式结果};
  \draw[study map arrow] (condition.north) -- (concept.south west) node[study map label,midway,left] {限定适用};
  \draw[study map arrow] (concept.south east) -- (derivation.north) node[study map label,midway,right] {用于推导};
  \draw[study map arrow] (condition.south) -- (result.west) node[study map label,midway,left] {条件约束};
  \draw[study map arrow] (derivation.south) -- (result.east) node[study map label,midway,right] {推导得到};
\end{tikzpicture}%
}
\end{center}

\begin{takeawaybox}[title={导图阅读指南}]
先确认术语定义及适用条件，再读推导与应用，最后核对关键结论成立的条件。实际生成时，把此段替换为依据资料撰写的关系解释，并回指对应正文。
\end{takeawaybox}
\StudyDivider

% 正文章节：概念正文必须先于图表；图表不得单独充当知识覆盖。
\section{章节标题}
\begin{center}
  {\small\textbf{说明：}本章覆盖……，核心概念包括……。重点掌握……，了解……即可。}
\end{center}
\vspace{0.3em}

\subsection{核心概念}
\StudyTerm{\PriorityMust}{术语名称（English Term）}{先给出准确、可独立复习的定义。}

\paragraph{为什么重要。}
说明它解决什么问题，以及它在本章知识结构中的位置。

\paragraph{条件与边界。}
列出成立条件、适用范围和不适用情形。

\paragraph{与相邻概念的关系。}
指出容易混淆的概念以及核心区别。

\paragraph{考试提示与常见误区。}
给出可操作的识别线索，不用图表替代概念解释。

\begin{definitionbox}[title={关键结论}]
内容框只放简短结论、记忆链条或复习提示；长推导和长案例留在普通正文中。
\end{definitionbox}

\subsection{推导与应用}
先说明目标、假设和变量，再展示必要的中间步骤：
\begin{align*}
  y &= a+b, \\
  z &= y+c = a+b+c.
\end{align*}

\begin{formulabox}[title={公式速记}]
\[
  z=a+b+c.
\]
\end{formulabox}

\subsection{必要的比较}
只有在表格能显著减少重复文字时才使用：
\begin{center}
\small
\begin{tabularx}{0.96\textwidth}{>{\bfseries}p{0.22\textwidth}XX}
\toprule
比较维度 & 理论 A & 理论 B \\
\midrule
核心含义 & …… & …… \\
适用条件 & …… & …… \\
考试识别 & …… & …… \\
\bottomrule
\end{tabularx}
\end{center}

\StudyDivider

\section{公式速查手册}
\begin{center}
  {\small\textbf{说明：}按主题分表，只汇总正文已经解释和推导过的关键公式；\PriorityMust\ 表示必须能默写。}
\end{center}

\subsection{主题一：基础公式}
\begin{FormulaSummaryTable}
  \FormulaSummaryRow{示例公式}{$z=a+b+c$}{\PriorityMust\ 必背；说明变量含义、成立条件与单位。}
  \FormulaSummaryRow{边界或变形式}{$a=z-b-c$}{只在需要反解参数时使用；不要省略符号条件。}
\end{FormulaSummaryTable}

\subsection{主题二：推断、检验或扩展公式}
\begin{FormulaSummaryTable}
  \FormulaSummaryRow{示例统计量}{$T=(\widehat\theta-\theta_0)/\operatorname{se}(\widehat\theta)$}{写明原假设、参考分布、自由度或大样本条件。}
\end{FormulaSummaryTable}

\section{背诵优先级速览}
\subsection{第一优先级 \PriorityMust（必背）}
\textbf{章节 A}：术语 1、术语 2；\textbf{核心公式}：公式 1。

\subsection{第二优先级 \PriorityImportant（重点）}
术语 3、术语 4。

\subsection{第三优先级 \PriorityKnow（了解）}
术语 5、拓展案例。

\section{全局易错点}
\begin{mistakebox}[title={考前最后检查}]
\begin{enumerate}[leftmargin=2em]
  \item 不要用图表标题代替概念定义。
  \item 不要把恒等式、假设、估计量和结论混为一谈。
  \item 检查公式的适用条件、符号和边界情况。
\end{enumerate}
\end{mistakebox}

\StudyDivider
\begin{center}
  {\large\bfseries —— 祝考试顺利 ——}
\end{center}

\end{document}
